\section{Related Work}
\label{sec:related}

\subsection{AI-Assisted Peer Review}

Recent work demonstrates AI's capability to generate structured paper reviews. D'Amour et al.~\cite{damour2023simulated} show that LLMs can produce reviews matching expert quality on 60\% of evaluation dimensions. Our prior work~\cite{dellalibera2026review} demonstrates an 8-stage lifecycle for AI-simulated reviews, achieving 32\% score improvement across 14 papers. However, these systems operate at a single tier: individual paper reviews without cross-portfolio synthesis.

Yuan et al.~\cite{yuan2024llm} explore using LLMs for meta-review generation, consolidating multiple reviews into a single summary. While this aggregates opinions, it does not identify cross-cutting patterns or escalate issues to higher organizational tiers. Our architecture extends this with explicit hierarchical structure and bidirectional flow.

\subsection{Hierarchical Feedback Systems}

Software engineering has long employed hierarchical review structures. Code review systems like Gerrit and GitHub use multi-level approval: individual reviewers, module owners, and repository administrators~\cite{rigby2013convergent}. However, these systems lack automated pattern detection and priority escalation rules.

Agile sprint retrospectives~\cite{schwaber2020scrum} follow a two-tier structure: team-level retrospectives feed into cross-team coordination meetings. Our architecture formalizes this with explicit priority classification (P1/P2/P3 → PP1/PP2/PP3 → B1/B2/B3) and bidirectional flow protocols.

\subsection{Priority Classification and Issue Triage}

Bug triage systems classify issues by severity (critical, major, minor) based on impact and frequency~\cite{anvik2006automating}. Our P1/P2/P3 classification adapts this model to review feedback, using reviewer consensus as the primary signal. The key innovation is \textbf{priority escalation}: a P3 issue at the paper tier can escalate to PP1 at the panel tier if it appears across multiple papers, indicating a systemic problem.

Enterprise governance frameworks like ITIL~\cite{axelos2019itil} use hierarchical incident management: local incidents escalate to major incidents when they affect multiple services. Our board tier (B1/B2/B3) applies this concept to research program coordination, escalating module-level patterns to program-level strategic decisions.

\subsection{Cross-Cutting Concern Detection}

Aspect-oriented programming~\cite{kiczales1997aspect} identifies cross-cutting concerns (logging, security) that span multiple modules. Our panel tier performs analogous pattern detection across papers, identifying methodological weaknesses or architectural inconsistencies that cut across the portfolio.

Mining software repositories research~\cite{hassan2008road} detects patterns in commit history and issue trackers. Our system applies similar techniques to review corpora, identifying repeated phrases and shared concerns across multiple review documents.

\subsection{Meta-Review and Crowdsourcing Aggregation}

Conference program committees employ meta-reviewers who synthesize individual reviews into consolidated assessments~\cite{shah2018challenges}. OpenReview~\cite{soergel2013openreview} provides infrastructure for meta-review generation, but relies on manual synthesis without automated pattern detection. Our panel and board tiers automate meta-review generation while preserving human oversight for high-stakes decisions.

Crowdsourcing research addresses quality control in multi-agent systems~\cite{bernstein2015crowdsourcing}. The Dawid-Skene model~\cite{dawid1979maximum} estimates worker quality via EM algorithms, weighting aggregated judgments by reliability. Expectation-Maximization approaches~\cite{raykar2010learning} infer ground truth from noisy labels. Our confidence-based deferral mechanism (Section~\ref{sec:oversight}) applies similar principles: low-confidence classifications trigger human review, preventing unreliable AI judgments from propagating.

Collective intelligence research~\cite{woolley2010evidence,malone2010handbook} studies when groups outperform individuals. Our emergent pattern detection validates this: 37-52\% of high-priority issues appear only at higher tiers, suggesting hierarchical synthesis produces insights no individual reviewer generates. However, collective intelligence also warns of groupthink and cascade effects~\cite{lorenz2011social}; our confidence scoring and human oversight mitigate these risks.

\subsection{Gaps in Existing Work}

No prior work combines hierarchical structure, automated pattern detection, priority escalation, and bidirectional flow in a single architecture for AI review systems. Existing systems either operate at a single tier or rely on manual coordination between tiers. Our contribution is a formal architecture with explicit protocols for upward issue escalation and downward directive propagation, validated on 14 papers across 3 modules.